\documentclass[10pt,a4paper]{amsart}
\usepackage[margin=19mm]{geometry}
\usepackage{amsmath,amssymb,mathtools}
\usepackage[scr=rsfs,cal=euler]{mathalfa}
\usepackage{xfrac}
\usepackage[unicode,hidelinks,bookmarks=false]{hyperref}
\newcommand{\ie}{i.e.\ }
\newcommand{\cf}{cf.\ }
\newcommand{\resp}{respectively\ }
\newcommand{\Subsection}{\subsection}
\providecommand{\nobreakdash}{\nobreak\hbox{-}}
\setlength{\emergencystretch}{2em}
\multlinegap0pt
\usepackage{dsfont}
\usepackage{thm-restate}
\usepackage{enumitem}
\newcommand{\N}{\mathds{N}}
\newcommand{\R}{\mathds{R}}
\newcommand{\Z}{\mathds{Z}}
\newcommand{\B}[1]{\{0,1\}^{#1}}
\newcommand{\fixed}[1]{\hat{ #1 }}
\newcommand{\lift}{l}
\newcommand{\sym}[1]{\mathcal{S}_{#1}}
\newcommand{\sprod}[2]{#1^\top #2}
\newcommand{\define}{\coloneqq}
\newcommand{\card}[1]{\left\lvert #1 \right\rvert}
\newcommand{\brackets}[1]{\left\{ #1 \right\}}
\newcommand{\naturalsto}[2][]{\left[ #2 \right]_{#1}}
\DeclareMathOperator{\conv}{conv}
\DeclareMathOperator{\supp}{supp}
\DeclareMathOperator{\argmax}{argmax}
\DeclareMathOperator{\argmin}{argmin}
\newtheorem{theorem}{Theorem}
\newtheorem{lemma}[theorem]{Lemma}
\newtheorem{remark}[theorem]{Remark}
\newtheorem{corollary}[theorem]{Corollary}
\newtheorem{proposition}[theorem]{Proposition}
\theoremstyle{definition}
\newtheorem{definition}[theorem]{Definition}
\newtheorem{example}[theorem]{Example}
\newtheorem{observation}[theorem]{Observation}
% The e-print carries, at source0.tex line 133, the commented-out author line reproduced verbatim
% in the next line. It is a template placeholder of the double-blind template and is not an author
% of the article; the authors of the article are the two named below.
% \author{Anonymous authors}
\title{A Framework for Handling and Exploiting Symmetry in Benders Decomposition}
\author{Christopher Hojny and C\'edric Roy}
\date{Source: arXiv:2511.22251v4 (CC BY 4.0); the authors' own native \LaTeX{} source, set here in the editorial AMS \texttt{amsart} wrapper (the source's own class is \texttt{article}, which is not redistributed).}
\begin{document}
\maketitle
\noindent\textit{Editorial scope.} The counted claim of this unit is the article's Lemma 12 (source label \texttt{rmk:zeta-lb}), the existence of a vector of $\zeta$-variables that dominates the group averages and satisfies both constraint families of the extended formulation, delivered together with the authors' complete original proof. Retained uncounted support, in source order, is the whole passage that builds the extended formulation inside Section~3: the two constraint families (9) and (10) that define the $\zeta$-variables, the generic extended-formulation cut (11), Theorem~9, the example witnessing that the extended formulation is strictly weaker, the restated Proposition~11, the paragraph that names the two intermediary lemmas, and, after the counted proof, Lemma~13 with its complete proof. Every retained passage is a verbatim copy of the authors' own source lines. The article's earlier sections, its abstract, its introduction and its later sections and appendices are not retained: those passages refer to the assignment-problem formulation and to further results that lie outside this source-local core, and their cross-references cannot resolve inside a unit of this size. Four cross-references that point outside the retained core are delivered as reference-only adaptations, namely as links to the article's own abstract page whose plain text is the number each prints in the published article: the two references to Proposition~8 on source lines 705 and 716 are carried as the plain number 8, and the two references to equation~(3) on source lines 715 and 725 are carried as the plain number (3). Nothing else is adapted: apart from those three declared reference-only adaptations every character of the delivered passages is the authors' own. Printed numbering: the authors' own theorem-environment declarations are carried unchanged, with the single shared counter of the theorem, lemma, remark, corollary, proposition, definition, example and observation environments, and the counter is advanced so that the delivered PDF prints Theorem~9, Proposition~11, Lemma~12 and Lemma~13 and the displayed equations (9), (10) and (11) with exactly the numbers they print in the published article. Class substitution: the source class is the standard \texttt{article} class with the authors' own packages; the publisher-independent wrapper is AMS \texttt{amsart} carrying the authors' own macros, the \texttt{thm-restate} package used by their restated Proposition~11, and the \texttt{dsfont} blackboard-bold font used by their own \texttt{R} and \texttt{N} macros. Reference handling: the article's own TeX contains no reference list and its bibliography exists only in the authors' file \texttt{bib.bib}, which is kept verbatim in the eprint directory and is not redistributed; no retained passage that would need a citation is included, no reference entry is transcribed and no reference entry is invented. No mathematics is written, repaired or re-derived here.
 The e-print additionally carries, commented out at source0.tex line 133, a template author line naming Anonymous authors; that line is not an author of the article, and it is reproduced verbatim as a comment in the editorial front matter so that the author record of the delivered unit is complete.
\setcounter{equation}{8}
\setcounter{theorem}{8}
\begin{align}
  \sum_{j \in G} x_j &\leq k - 1 + \left(|G| - (k - 1)\right)\zeta^k_G, && G \in \mathcal{G},\; k \in [|G|],\label{eq:zeta-lowerbound}\\
  k \cdot \zeta^k_G &\leq \sum_{j \in G} x_j, && G \in \mathcal{G},\; k \in [|G|].\label{eq:zeta-upperbound}
\end{align}
%
The variables~$\zeta^1_{G}, \dots, \zeta^{|G|}_{G}$ can thus be considered a non-increasing reordering of~$\brackets{x_{j}: j \in G}$.
By denoting the~$k$-th largest value in~$\{\alpha_j : j \in G\}$ by~$\alpha^k_{G}$, the most violated cut constructed in the proof of Proposition~\href{https://arxiv.org/abs/2511.22251v4}{8} can be phrased in terms of the~$\zeta$-variables as
%
\begin{equation}
  \sum_{G\in \mathcal{G}} \sum_{k = 1}^{|G|} \alpha^k_{G} \zeta^k_{G} +\beta \leq z. \label{eq:generic-EF-cut}
\end{equation}
%
We thus immediately obtain the following theorem.
%
\begin{theorem}
  \label{thm:EF-works}

  Assume all~$x$-variables in~\href{https://arxiv.org/abs/2511.22251v4}{(3)} are binary.
  Let~$(\alpha,\beta) \in \mathcal{I}$ and symmetry group~$\Pi$ adhere to the conditions in Proposition~\href{https://arxiv.org/abs/2511.22251v4}{8}.
  Then, $(\fixed{x},\fixed{z}) \in \B{n} \times \R$ satisfies all equivalent inequalities in~$\{(\pi^{-1}(\alpha),\beta) : \pi \in \Pi\}$ if and only if~$(\fixed{x},\fixed{\zeta},\fixed{z})$ in the extended model satisfies~\eqref{eq:generic-EF-cut}.
\end{theorem}

\setcounter{theorem}{9}
%

We now have two approaches to benefit from symmetries when generating Benders cuts: the symmetric cut pool and the extended formulation.
It is unclear though which of these two formulations is stronger.
In particular, the extended model is only an extended IP formulation, but we have not settled whether it is also an extended formulation of the LP relaxation of the cut pool model.
In the following, we answer this question rigorously.
While Theorem~\ref{thm:EF-works} shows that the separation of symmetric Benders cuts can be avoided, the extended model~\eqref{eq:zeta-lowerbound}--\eqref{eq:generic-EF-cut} has a weaker LP relaxation than~\href{https://arxiv.org/abs/2511.22251v4}{(3)} in general.
Indeed, Proposition~\ref{lemma:EF-is-weaker} shows that every feasible~$(x, z)\in [0,1]^{n}\times\R$ induces a feasible~$(x, \zeta, z)\in[0,1]^{2n}\times\R$.
The converse is however not true.
The next example shows that there exist some~$(x,\zeta, z)$ with projection to~$(x,z)$ not satisfying~$\sprod{\alpha}{x} + \beta \leq z$.

\begin{example}
  Consider a problem with three variables $x$, $y_1$, and $y_2$ for which the conflict cuts $x + y_1 \leq 1$ and $x + y_2 \leq 1$ can be derived.
  Assuming that $y_1$ and $y_2$ are symmetric, the two conflict cuts belong to a symmetric cut pool, which can be represented by $\zeta^1_x + \zeta^1_y \leq 1$ in the extended model.
  Then, for every $\varepsilon \in [0,\frac{1}{2}]$, the solution $(x,y_1,y_2,\zeta^1_x,\zeta^1_y,\zeta^2_y) = (1-\varepsilon, 2\varepsilon, 0, 1-\varepsilon,\varepsilon,\varepsilon)$ is feasible for the LP relaxation of the extended model.
  However, solution $(x,y_1,y_2) = (1-\varepsilon,2\varepsilon,0)$ is not feasible for the cut pool model whenever~$\varepsilon > 0$.
\end{example}

\begin{restatable}{proposition}{lemmaEFisWeaker}
  \label{lemma:EF-is-weaker}
  Let~$(x,z)\in[0, 1]^{n}\times\R$, $\alpha\in\R^n_{\geq 0}$ satisfy~$\sum_{j\in C} \alpha_j \pi(x)_j +\beta \leq z$ for all~$\pi\in\Pi$.
  Then there exists a~$\zeta$ such that~$(x,\zeta, z)$ is feasible in the relaxation of the extended formulation.
\end{restatable}

The proof of Proposition~\ref{lemma:EF-is-weaker} requires two intermediary results.
Namely, we first show that we can construct~$\zeta$ such that its entries are upper-bounded by the average of~$x$ within the corresponding group~$G$ (Lemma~\ref{rmk:zeta-lb}).
We can then show that such a~$\zeta$ will satisfy constraint~\eqref{eq:generic-EF-cut}.
Trivially, the sum of the first~$k$ largest element of~$x$ is always greater than~$k$ times the average.
We show that this property still holds in a weighted sum when the weights are non-increasing (Lemma~\ref{rmk:weighted-average-dom}).

\begin{lemma}
  \label{rmk:zeta-lb}
  Let~$x\in[0, 1]^n$. 
  There exists~$\zeta\in [0, 1]^n$ satisfying~\eqref{eq:zeta-lowerbound} and~\eqref{eq:zeta-upperbound} such that~$\zeta^k_G \leq \frac{x(G)}{|G|}$ for all~$G\in\mathcal{G}$ and~$k\leq|G|$.
\end{lemma}

\begin{proof}
  For ease of notation, let~$x(S)$ denote~$\sum_{j\in S} x_j$ for any set~$S\subseteq\naturalsto[]{n}$.
  Combining Equations~\eqref{eq:zeta-lowerbound} and~\eqref{eq:zeta-upperbound}, we get that
  \begin{equation*}
    \frac{x(G) -k + 1}{|G| - k + 1} \leq \zeta^k_G \leq \frac{x(G)}{k}
  \end{equation*}
  which is a non-empty interval for any~$1\leq k\leq |G|$ and~$x\in [0, 1]^n$.
  Indeed,
  \begin{equation*}
    \frac{x(G)}{k} - \frac{x(G) - k + 1}{|G|-k+1} = \frac{x(G)(|G| -2k+1) + k(k-1)}{k(|G|-k+1)}.
  \end{equation*}
  Observe that the numerator can be rewritten as a function of~$k$ or of~$x(G)$ in the following:
  \begin{equation*}
    x(G)\left(|G| + 1 - 2k\right) + k(k-1) = k^2 - (2x(G)+1)k + x(G)(|G|+1).
  \end{equation*}
  The minimum with respect to~$x(G)$ is either at~$x(G) = 0$ or~$|G|$, depending on whether~$k \geq \frac{|G|+1}{2}$. 
  Reciprocally, the minimum with respect to $k$ is when~$2k - (2x(G) + 1) = 0$, or equivalently when~$k$ is close to~$\frac{1}{2} + x(G)$.
  As~$k\in \{1, \dots, |G|\}$, the global minimum of the numerator is either when~$k = 1, x(G) = 0$ or~$k=x(G)=|G|$.
  In either case, the numerator becomes~$0$, meaning that the interval for~$\zeta^k_G$ is never empty.
  Let~$\hat{\zeta}^k_G$, $k\leq |G|$ be the unique~$\zeta$-values that satisfy~\eqref{eq:zeta-lowerbound} with equalities.
  We can observe that for all~$G\in\mathcal{G}$, $\hat{\zeta}^1_G = \frac{x(G)}{|G|}$ and for all~$k\leq|G|-1$ we have
  \begin{align*}
    \hat{\zeta}^k_G - \hat{\zeta}^{k+1}_G = \frac{x(G) - k + 1}{|G| - k + 1} - \frac{x(G) - k}{|G| - k} = \frac{|G| - x(G)}{(|G|-k+1)(|G|-k)} \geq 0.
  \end{align*}
  Observe that~$\hat{\zeta}^k_G$ satisfying~\eqref{eq:zeta-lowerbound} with equality does not always guarantee~$\hat{\zeta}^k_G \geq 0$.
  Nevertheless, as the upper bound derived from~\eqref{eq:zeta-upperbound} is always greater or equal to zero,  $\zeta^k_G = \max\brackets{\hat{\zeta}^k_G, 0}$ ensures that~$\frac{x(G)}{|G|} \geq \zeta^1_G \geq \zeta^2_G \geq \dots \geq \zeta^{|G|}_G$ as claimed.
\end{proof}


\begin{lemma}
  \label{rmk:weighted-average-dom}
  Let~$\alpha, x\in \R^n$ such that~$\alpha_1 \geq \dots \geq \alpha_{n} \geq 0$ and~$x_1 \geq \dots \geq x_{n} \geq 0$.
  Denoting the average of~$x$ by~$A_v(x)\define \sum^n_{k=1} x_i/n$, then $\sum_{k=1}^K \alpha_k A_v(x) \leq \sum_{k=1}^K \alpha_k x_k$ for any~$K\leq n$.
\end{lemma}
\begin{proof}
  Let~$f(K)\coloneq \sum_{k=1}^K \alpha_k (x_k - A_v(x))$. 
  We prove that~$f(K)\geq 0$ for any~$K\leq n$.
  Observe that~$f$ behaves like a concave function as~$\alpha_k$ and~$(x_k - A_v(x))$ both decrease when~$k$ increases and~$\alpha_k \geq 0$.
  We then only need to check that~$f(1) \geq 0$ and~$f(n)\geq 0$.
  Since~$x_1 \geq A_v(x)$ and $\alpha_1 \geq 0$, we conclude $f(1) \geq 0$.
  Let~$g(k)\coloneqq \sum_{i=1}^k (x_i - A_v(x))$.
  It is easy to verify that~$g(k)\geq 0$ for the same arguments:~$g$ also behaves like a  concave function, $g(1)\geq 0$, and~$g(n) = 0$.
  Notice that since~$g(k) - g(k-1) = x_k - A_v(x)$ for all~$k\in \brackets{2, ..., n}$, we have that
  \begin{align*}
    f(n)
    &= \sum_{k = 1}^n \alpha_k(x_k - A_v(x)) = \sum_{k = 2}^n \alpha_k(x_k - A_v(x)) + \alpha_1 (x_1 - A_v(x)) \\
    &= \sum_{k=2}^{n} \alpha_k (g(k) - g(k-1)) + \alpha_1 g(1) = \sum_{k=2}^{n} \alpha_k g(k) - \sum_{k=2}^{n} \alpha_k g(k-1) + \alpha_1 g(1) \\
    &= \alpha_n g(n) + \sum_{k=2}^{n-1} \alpha_k g(k) - \sum_{k=1}^{n-1} \alpha_{(k+1)} g(k) + \alpha_1 g(1) \\
    &= \alpha_n \underbrace{g(n)}_{=0} + \sum_{k = 1}^{n-1} (\underbrace{\alpha_k - \alpha_{k+1}}_{\geq 0}) \underbrace{g(k)}_{\geq 0} 
  \end{align*}
  and thus~$f(n) \geq 0$.
\end{proof}

\end{document}
